\section{Tercera iteración de la Fase 5: Planificar el desarrollo del generador}
\label{sec:F5I3}

Tras el análisis y la definición del formato, el código que transforma las listas finales del
editor en ese fichero, su integración en el editor con el botón
«Generar GWI+GWC», las estadísticas y sus pruebas.

\subsection{El módulo \texttt{BaseCartucho.js}}
\label{sec:F5I3Modulo}

Todo lo relativo a los cartuchos vive ahora en un módulo javascript
nuevo, \texttt{BaseCartucho.js}, cargado tras el descomponedor, con
una única responsabilidad: la \textbf{lectura y la grabación} de
los tres formatos de cartucho del proyecto:

\begin{itemize}
	\item \textbf{GWB}: Cartucho Base, con los gráficos que se usarán en el editor gráfico y los nombres de los botones.
	
	\item \textbf{GWI}: El Cartucho Intermedio con el formato de autoría, la definciión de los segmentos y las imágenes PNG.
	
	\item \textbf{GWC}: El cartucho Compacto con el formato de ejecución.
\end{itemize}

Hasta ahora esa lógica
estaba repartida — el guardado del GWI en el descomponedor, la
carga del GWB y del GWI en el director — y el GWC iba a nacer
suelto; con el módulo, el reparto queda limpio: el editor edita,
el motor juega, \texttt{BaseCartucho.js} lee y escribe ficheros de los tres cartuchos que manejamos. Cada función está separada para cada tipo de cartucho y son todas independientes.

El módulo no conoce al editor ni el DOM: recibe la estructura que existe en memoria y produce ficheros, o bien recibe ficheros y reconstruye la estructura.

Las funciones de lectura y validación del GWC quedan en una \textbf{sección delimitada y autocontenida} del fichero, que no depende del GWB ni del ZIP: servirá, con los cambios oportunos, como base del lector del firmware del ESP32.

El generador no usa el azar en ninguna parte: recorre las listas
siempre en el mismo orden y escribe siempre lo mismo. La generación
es \textbf{determinista}: el mismo GWI produce el mismo GWC byte a
byte, con la misma suma de verificación\footnote{En el apartado \ref{sec:F5I14CRC} se reemplaza la suma de verificación simple por un CRC-16} — el equivalente del
compilador reproducible, y la base de su prueba de regresión.

El descarte también vive aquí: los identificadores legibles, los
fantasmas tipo 99, los enlaces inversos calculados y toda la sección
de autoría no viajan; el generador los descarta al vuelo, un
segmento que nadie enciende ni pinta no existe en el GWC.\\\

\begin{table}[H]
	\centering
	\begin{tabularx}{\textwidth}{@{}l l X@{}}
		\toprule
		\textbf{Formato} & \textbf{Función} & \textbf{Descripción} \\
		\midrule
		
		\texttt{GWB} & \texttt{leerGWB} & Carga el cartucho de gráficos de
		la máquina tras el setup: sus imágenes y sus botones \\
		\midrule
		
		\texttt{GWI} & \texttt{leerGWI} & Carga el cartucho de autoría: el
		JSON y las imágenes PNG del ZIP, y reconstruye las listas en
		memoria igual que desde el editor \\
		
		& \texttt{guardarGWI} & Graba el cartucho en ZIP con el JSON+imágenes \\
		\midrule
		
		\texttt{GWC} & \texttt{leerGWC(bytes)} & La puerta de entrada de la lectura:
		recorre la cabecera y los bloques, valida la suma de
		verificación\footnote{En el apartado \ref{sec:F5I14CRC} se reemplaza la suma de verificación simple por un CRC-16} y reconstruye las listas en memoria con las
		mismas formas que produce el editor \\
		
		& \texttt{guardarGWC(listas)} & Valida los topes, escala, ensambla los bloques, escribe la cabecera con la suma de verificación\footnote{En el apartado \ref{sec:F5I14CRC} se reemplaza la suma de verificación simple por un CRC-16} y devuelve los bytes junto a la estadística de tamaños \\
		
		& \texttt{escalar(sprite, pantalla)} & Convierte un sprite en su
		máscara de 1 bit escalada a la pantalla declarada, con el
		muestreo por vecino más cercano y escala por eje del
		simulador: lo que viaja es lo que el simulador dibuja \\
		
		& \texttt{indiceDe(id)} & Traduce los \texttt{seg\_NNN} del GWI a
		índices del GWC, con el \textbf{nulo} para los enlaces
		ausentes; los campos \emph{salvadoContador} y
		\emph{muerteContador} viajan con el contador correcto o nulos \\
		
		& \texttt{bloque(tipo, contenido)} & Escribe la estructura común
		de todos los bloques: tipo, tamaño del contenido de 4 bytes y
		contenido \\
		
		\midrule
		\texttt{Común} & \texttt{validarCartucho(listas)}\footnotemark & Comprueba que se puede generar correctamente los cartuchos con la información definida, si hay errores bloquea la generación.\\	
		
		& \texttt{validarTopes(listas)} & Comprueba que sprites, segmentos
		y secuencias no superan los 255 elementos direccionables con
		índices de 1 byte: si alguno pasa, el error es bloqueante y no
		se genera nada \\
		
		& \texttt{estadistica(listas, bytes)} & Las cifras acordadas: la
		ocupación de los sprites sin deduplicación frente a la final, y
		el tamaño del GWI completo frente al GWC \\
		
		\bottomrule
	\end{tabularx}
	\caption{Las funciones de \texttt{BaseCartucho.js} por formato de cartucho.}
	\label{tab:funcionesGenerador}
\end{table}
% El texto de la nota se coloca justo aquí, fuera del entorno de la tabla
\footnotetext{Se decidió en la iteración 5 hacer este cambio, se incluye aquí por completitud}


\subsection{El botón «Generar GWI+GWC»}
\label{sec:F5I3Boton}

La generación se lanza desde el editor gráfico, que tiene toda la información, por lo que se añade un nuevo botón junto al de «Verificar y Guardar» que genera el GWI (este se mantiene para poder dejar un cartucho a medias y seguir en otro momento), así que el GWI se puede guardar en cualquier momento sin más verificación que la suya propia. El nuevo botón encadena los dos procesos:

\begin{enumerate}
	\item \textbf{Verificación y guardado del GWI}, idéntico al flujo actual: la verificación completa del diseñador, bloqueante, y el ZIP con el JSON y las imágenes. Se marcarça como $\text{Versión} = 0$ para indicar que es un fichero auxiliar que no ha generado el GWC.
	
	\item \textbf{Generación del GWC} sobre las mismas listas en 	memoria que acaban de producir el GWI, con su descarga (\texttt{cartucho\_gwc.gwc}).
\end{enumerate}

La razón del encadenamiento es por comodidad y atomicidad:
\begin{itemize}
	\item La verificación se lanza siempre antes de generar el GWI, al generarlos a la vez no hay que volver a validar los datos.
	\item Si la validación falla, no se genera ninguno de los dos; si tiene éxito, ambos salen de la misma información de la memoria y son \textbf{siempre copia uno del otro}.
	\item No existe la situación de un GWC huérfano de un GWI más nuevo. Los avisos de la barra de estado acompañan cada paso («Guardando GWI…», «Generando GWC…»), y cualquier fallo es visible con su mensaje, como en todo el proyecto.
\end{itemize}

Al limitar los punteros a un byte, es necesario verificar que no existen mas de 256 segmentos. Esto se efectuará en tres lugares: 

\begin{itemize}
	\item Tras descomponer los segmentos dará un aviso, ya que pueden unirse varios y reducir la cantidad.
	\item Al guardar el GWI dará otro aviso, todavía se puede seguir trabajando en los segmentos.
	\item Al pulsar «Generar GWI+GWC» es ya un error \textbf{bloqueante} pues no se puede generar el GWC con errores. 
\end{itemize}

De esta manera se avisa durante el trabajo, pero solos e bloquea cuando e imprescindible, lo que mejora el flujo de trabajo al usuario, es muy frustrante que tras definirlo todo el sistema de un error al final de todo, cuando puede avisar antes.

\subsection{Simulación}
\label{sec:F5I3Simulacion}

El generador es puro (lista en memoria $\rightarrow$ bytes, sin
DOM), así que la batería de pruebas lo ejercita entero con listas
sintéticas y con el GWI real del Parachute:

\begin{itemize}
	\item \textbf{Determinismo:} dos generaciones del mismo GWI
	producen bytes idénticos y la misma suma de verificación\footnote{En el apartado \ref{sec:F5I14CRC} se reemplaza la suma de verificación simple por un CRC-16}.
	\item \textbf{Estructura:} los conteos de la cabecera cuadran con
	las tablas; toda dirección apunta dentro del fichero; la suma de
	verificación\footnote{En el apartado \ref{sec:F5I14CRC} se reemplaza la suma de verificación simple por un CRC-16} recomputada coincide; los punteros de los bloques
	opcionales ausentes son nulos.
	\item \textbf{Contenido:} los enlaces del registro se resuelven
	contra los índices reales; ningún enlace válido apunta al nulo ni
	al revés; \emph{salvadoContador} y \emph{muerteContador} viajan
	correctos o nulos; no viajan fantasmas tipo 99 ni absorbidos.
	\item \textbf{Escalado:} posiciones y máscaras coinciden a píxel
	con lo que el simulador pinta al vuelo con la misma pantalla
	declarada.
	\item \textbf{Topes:} un GWI con más de 255 segmentos produce el
	error bloqueante y ningún fichero.
	\item \textbf{Fondo:} la versión 1 no escribe el bloque de fondo;
	si el GWI declara color de fondo, viaja en el Bloque 0.
\end{itemize}

Y en navegador real, con el cartucho del Parachute: el botón produce
ambos ficheros seguidos; un cartucho con un error de verificación no
produce ninguno; las cifras de la estadística se presentan en la
barra; la inspección del GWC (volcado hexadecimal) muestra el
identificador «GWC», su versión y el recorrido de bloques en orden.

\subsection{La estadística de tamaños}
\label{sec:F5I3Estadistica}

Al terminar, la generación presenta en la barra de estado y en
detalle bajo GW\_DEBUG las cifras acordadas:

\begin{itemize}
	\item \textbf{Ahorro de la deduplicación:} la ocupación que
	tendrían todos los sprites sin deduplicación (una máscara por
	aparición) frente a la ocupación final de los sprites distintos en
	el GWC, con la cantidad ahorrada.
	\item \textbf{Ahorro total:} el tamaño del cartucho GWI completo
	(el ZIP con el JSON y todos sus gráficos) frente al tamaño final
	del GWC, con la cantidad total ahorrada.
\end{itemize}

El Parachute real es el caso de estudio: sus números concretos
—recogidos en la memoria al cierre de la fase— ilustran la
distancia entre un formato de autoría (legible, editable, con
historial e imágenes PNG) y uno de ejecución (binario, mínimo,
listo para el microcontrolador).

\subsection{Carga en el simulador y Verificación}
\label{sec:F5I3Carga}

Se modificará el generador gráfico, de forma que permita cargar un
cartucho en formato GWC con \texttt{leerGWC}: de esta manera se
puede verificar que su contenido es válido, que visualmente carga
lo mismo que un GWI, que mantenga las propiedades correctas en los
segmentos, y que el simulador ejecute el mismo juego.

Además esto permite una verificación bastante sencilla con el ciclo
de ida y vuelta:

\begin{enumerate}
	\item Se carga un GWI.
	\item Sin tocar nada se genera el GWI+GWC.
	\item Se carga el GWC generado.
	\item Sin tocar nada se genera de nuevo el GWI+GWC.
\end{enumerate}

Los dos GWC —el del punto 2 y el del punto 4— deben ser
\textbf{idénticos byte a byte}, lo que es muy sencillo de validar
con comandos estándar del sistema. El GWI regenerado es
\textbf{equivalente en su contenido de juego}: la autoría y los
gráficos originales no viajan en el GWC, así que los GWI no pueden
ser idénticos, pero las listas que producen deben ser las mismas.

El resto de pruebas de la carga (estructura, contenido, escalado,
topes y fondo) son las de la batería de la sección anterior; el
cierre de la fase añadirá la batería de pruebas con su resultado,
como en las fases anteriores, y las cifras reales del Parachute
para la memoria.